\section{A single diagonal sequence for every bounded-test energy sublevel}
All assertions in this section are deductions on the original actual
carrier and class, under the analytic foundations reconstructed above.
They impose no restriction on the support of the prescribed measure.
They do not recognize a new function as a CLD psh weight.

Fix any eligible $\mu$, write $B=D(\mu)$, and use the common maximizing
tail bound $C=C_\mu$. Put $\delta_n=2^{-4n}$, $\beta_n=2^{-n}$ and
$R_j=\sum_{n\ge j}\delta_n$. For $L\ge0$ define, on Borel sets,
\[
 \operatorname{Cap}_L(A)=\sup_{w\in\C_L}\MA(w)(A).
\]
This is an actual analytic-space set function, not a substitute MA
operator. It is monotone and countably subadditive, by those properties
of each of its defining positive measures; $\operatorname{Cap}_L(X)=2$.
No regularity or recognition theorem for this capacity is assumed.

\begin{theorem}\label{new:diagonal}
There is one normalized actual maximizing sequence $a_n\in\C_C$,
with $\varepsilon_n=B-F_\mu(a_n)\to0$, such that, for every integer
$k\ge1$ and $n\ge k$,
\begin{equation}\label{new:increment}
 \sup_{w\in\C_{C+k}}\int|a_{n+1}-a_n|\MA(w)\le\delta_n.
\end{equation}
The same bound holds with the canonical weak limit $\nu$ in place of
$\MA(w)$. There is a decreasing sequence of ordinary open sets $O_m$
such that, for $m\ge k$,
\begin{equation}\label{new:capacity}
 \operatorname{Cap}_{C+k}(O_m)\le P_m:={8\over7}2^{-3m}.
\end{equation}
The Borel set $N=\bigcap_mO_m$ is null for every bounded MA measure,
for every bounded mixed measure, and for every weak limit of measures
from any fixed $\C_L$. It contains no point of full graded dimension
two. On $X\setminus O_m$, $a_n$ converges uniformly to a finite
continuous function, and
\[
 \sup_{X\setminus O_m}|a_n-u|\le 2^{1-n}\quad(n\ge m),
 \qquad u:=\limsup_n a_n.
\]
Here $u$ is an actual Borel function on all $X$, valued in $[-\infty,0]$.
It is finite outside $N$ and agrees with $\ell$ at every full point.
Its continuity is only on the specified closed complements. Global
USC, psh positivity and generated membership are not asserted.
\end{theorem}
\begin{proof}
Let $L_n=C+n+1$. Since $T_{L_n}(q)\to0$ as $q\downarrow0$, choose
$r_n>0$ with $r_n\le2^{-n}$ and $T_{L_n}(r_n)\le\delta_n$.
Choose actual normalized competitors with
\[
 \varepsilon_n\le\min\{1,2^{-n},r_{n-1}/8,r_n/8\},
\]
omitting $r_0$ at $n=1$. Such choices exist from the real supremum;
coercivity puts them in $\C_C$. Alternatively they can be selected
from any given maximizing net using increasing indices and directedness.
There is no assertion of countable cofinality. Equation \eqref{deficit}
gives $I(a_n,a_{n+1})\le r_n$. For $n\ge k$, the two endpoints and
every $w\in\C_{C+k}$ belong to $\C_{L_n}$, so \eqref{third} proves
\eqref{new:increment}. Passing the inequality for the fixed continuous
function $|a_{n+1}-a_n|$ to $\MA(a_i)\rightharpoonup\nu$ proves its
$\nu$ version.

Set $G_n=\{|a_{n+1}-a_n|>\beta_n\}$ and
$O_m=\bigcup_{n\ge m}G_n$. These are ordinary open sets. Markov's
inequality and countable subadditivity give, for $m\ge k$,
\[
 \operatorname{Cap}_{C+k}(O_m)
 \le\sum_{n\ge m}\delta_n/\beta_n
 =\sum_{n\ge m}2^{-3n}=P_m.
\]
Consequently $N$ is null for every $\MA(w)$: normalize $w$ and choose
$k$ with $-E(w)\le C+k$. Since
$B(p,q)\le2\MA((p+q)/2)$ and normalizing the midpoint preserves its
curvature and increases its energy, the same conclusion holds for
all bounded mixed measures.

We spell out the weak-limit assertion without a sequential compactness
assumption. If $\lambda_i\rightharpoonup\lambda$, with
$\lambda_i=\MA(w_i)$ and $w_i\in\C_L$, then for any ordinary open
$O$,
\[
 \lambda(O)\le\operatorname{Cap}_L(O).
\]
Indeed, compact-Hausdorff regularity and Urysohn's lemma express
$\lambda(O)$ as the supremum of $\int h\,d\lambda$ over continuous
$0\le h\le1$ supported in $O$. Each such integral is the limit of
$\int h\,d\lambda_i\le\operatorname{Cap}_L(O)$. Taking the supremum
proves the assertion. Choose $k$ with $L\le C+k$ and apply it to
$O_m$ for $m\ge k$; then $\lambda(N)=0$. In particular $\nu(N)=0$.

For a full point $x$, an actual atom
$\MA(w_x)\ge\alpha_x\delta_x$ has its normalized test in some
$\C_{C+k}$. If $x\in N$, monotonicity would give
$\alpha_x\le\operatorname{Cap}_{C+k}(O_m)\le P_m$ for every
$m\ge k$, a contradiction. Thus no full point lies in $N$.
If $y\notin O_m$, every increment from index $m$ on has magnitude
at most $\beta_n$. The telescoping bound
$|a_q(y)-a_n(y)|\le\sum_{l=n}^{q-1}\beta_l\le2^{1-n}$,
uniform for $y\in X\setminus O_m$, proves uniform convergence.
This closed subset is compact, and $a_m$ is bounded there, so the limit
is finite and continuous there. The pointwise limsup is Borel and
equals that limit. At all full points it equals the previously proved
finite, family-independent tuple $\ell$. Only countable measure
operations have occurred.
\end{proof}

\section{Integral tails and actual finite maxima: a stronger bridge}
Write $h_j=\sum_{n\ge j}|a_{n+1}-a_n|$, a nonnegative Borel function,
and $M_{j,q}=\max(a_j,\ldots,a_q)$, $S_j=\sup_{n\ge j}a_n$.
All $M_{j,q}$ are in the entire actual class $\C_C$. The $S_j$ are
bounded finite LSC functions, $a_j\le S_j\le0$, with $S_j\downarrow u$.
They have no proved positivity class.

\begin{theorem}\label{new:tails}
For $j\ge k$, uniformly over $w\in\C_{C+k}$ and $q\ge j$,
\begin{align}
 \int h_j\MA(w)&\le R_j,\label{new:variation}\\
 \int|a_j-u|\MA(w),\quad\int|M_{j,q}-u|\MA(w),\quad
 \int(S_j-u)\MA(w),\quad\int(S_j-a_j)\MA(w)&\le R_j.
 \label{new:alltails}
\end{align}
The same conclusions hold for $\nu$. The integrals are finite; values
on the common null set may be assigned arbitrarily when necessary.
Furthermore,
\begin{equation}\label{new:maximum}
 0\le E(M_{j,q})-E(a_j)\le R_j,
 \qquad 0\le I(a_j,M_{j,q})\le R_j.
\end{equation}
Uniformly for every $q\ge j$ as $j\to\infty$, the actual bounded
MA measures of $M_{j,q}$ converge weakly to $\nu$, and their energies,
reference moments, self mixed moments and diagonal MA moments tend to
the same limits as for $a_j$. The entire two-index family $M_{j,q}$
is maximizing for the full dual of $\nu$ in this uniform sense.
\end{theorem}
\begin{proof}
Countable Tonelli applied to \eqref{new:increment} proves
\eqref{new:variation}. It also proves the $\nu$ bound from its
increment bounds. Outside $N$, the sequence has finite total variation:
after some index all its increments are at most $\beta_n$. For
$n\ge j$, $|a_n-u|\le h_n\le h_j$, and
$|a_n-a_j|\le h_j$. Taking finite maxima and the supremum proves
$|M_{j,q}-u|\le h_j$, $0\le S_j-u\le h_j$, and
$0\le S_j-a_j\le h_j$. Since all the measures in question annihilate
$N$, integration proves \eqref{new:alltails}. In particular $u$ is
integrable against each bounded MA measure, each bounded mixed measure
and $\nu$, and $\int u\,d\nu=p$ by the already proved
$\int a_j\,d\nu\to p$.

Apply tangent concavity at $a_j$:
\[
 0\le E(M_{j,q})-E(a_j)
 \le\int(M_{j,q}-a_j)\MA(a_j)\le R_j.
\]
Green and actual order give
\[
 0\le I(a_j,M_{j,q})
 =\int(M_{j,q}-a_j)[\MA(a_j)-\MA(M_{j,q})]
 \le\int(M_{j,q}-a_j)\MA(a_j)\le R_j.
\]
All tests are on $\C_C\subset\C_{C+1}$, so no diagonal restriction
on $j$ is needed for these assertions. For a fixed normalized test $w$,
\eqref{weaktests} bounds its MA-measure discrepancy by
$\sqrt{12D R_j}$, where $D=\max(C,-E(w),0)$.
Uniform density and common mass extend the conclusion to every
continuous test. The bounds \eqref{momentCauchy}, applied to
$a_j,M_{j,q}$, prove the stated uniform moment limits; in particular
the diagonal error is at most $H_C(R_j)+\sqrt{12CR_j}$.

Finally let $d=E(M_{j,q})-E(a_j)$ and
$z=\int(M_{j,q}-a_j)\,d\nu$. Both belong to $[0,R_j]$.
Thus $|F_\nu(M_{j,q})-F_\nu(a_j)|=|d-z|\le R_j$.
Since $a_j$ also maximizes for $\nu$, the full family does as claimed.
No energy or operator has been assigned to $S_j$ by this proof.
\end{proof}

For completeness, $L_j=\sup_{q\ge j}E(M_{j,q})$ exists,
$E(a_j)\le L_j\le E(a_j)+R_j$, and $L_j\downarrow e$.
The energies increase with $q$. Comparing
$M_{j+1,q}\le M_{j,q}$, and also
$E(M_{j,j})\le E(M_{j,j+1})$, proves $L_{j+1}\le L_j$.
The squeeze by $E(a_j)$ and $R_j$ proves its limit. This independently
recovers the whole finite-tail-maximum working-note deduction and
strengthens it by \eqref{new:alltails} and \eqref{new:maximum}.

\section{An exact prescribed-measure absolute-integral defect}
\begin{theorem}\label{new:Lone}
For the same completely unrestricted prescribed $\mu$, the actual
Borel limsup $u$ is $\mu$-integrable, with possible $-\infty$ values
only on a $\mu$-null set. Put
\[
 s=\int u\,d\mu,\qquad r=e-B,\qquad J=s-r.
\]
Then $0\le J\le B-e$, and
\begin{align}
 \int(S_j-u)\,d\mu&\longrightarrow0,\label{new:orderedmu}\\
 \int(a_j-u)_+\,d\mu&\longrightarrow0,\label{new:onesidemu}\\
 \int|a_j-u|\,d\mu&\longrightarrow J.\label{new:exactLone}
\end{align}
Moreover the exact uniform maximality loss is
\begin{equation}\label{new:lostmax}
 \lim_{j\to\infty}\sup_{q\ge j}[B-F_\mu(M_{j,q})]=J,
 \quad
 \lim_{j\to\infty}\lim_{q\to\infty}F_\mu(M_{j,q})=B-J.
\end{equation}
Thus $J=0$ is equivalent both to $a_j\to u$ in $L^1(\mu)$ and to
uniform prescribed-measure maximality of all finite tail maxima.
These are exact deductions for the selected fast sequence. We do not
prove $J=0$, do not assert it is independent of the fast sequence,
and do not identify it with $K=B-D(\nu)$.
\end{theorem}
\begin{proof}
For each actual normalized $M_{j,q}$, the full dual gives
\[
 \int(-M_{j,q})\,d\mu\le B-E(M_{j,q}).
\]
At fixed $j$, $M_{j,q}\uparrow S_j$ and $|M_{j,q}|\le-a_j$,
a finite continuous function. Dominated convergence, or countable
monotone convergence of $M_{j,q}-a_j$, gives
\[
 \int(-S_j)\,d\mu\le B-L_j.
\]
Now $-S_j\uparrow-u$. Ordinary sequential monotone convergence and
$L_j\downarrow e$ yield $\int(-u)\,d\mu\le B-e<\infty$.
This proves integrability and \eqref{new:orderedmu}, using
$S_j\ge u$. Since $S_j\ge a_j$, passage to the limits in
$\int S_j\,d\mu\ge\int a_j\,d\mu$ gives $s\ge r$.
Since $u\le0$, also $s\le0$, proving the bounds on $J$.

Pointwise $(a_j-u)_+\le S_j-u$ outside the null pole set, so
\eqref{new:orderedmu} proves \eqref{new:onesidemu}. The elementary
identity for finite real values
\[
 |a_j-u|=u-a_j+2(a_j-u)_+
\]
holds almost everywhere and has finite integrals. Since
$\int a_j\,d\mu\to r$, it proves \eqref{new:exactLone}.
This proves an absolute-integral limit without inferring $L^1$
convergence from signed integrals.

Let $g_j=\int(S_j-a_j)\,d\mu$. Then $g_j\to s-r=J$.
For $q\ge j$ put $z_{j,q}=\int(M_{j,q}-a_j)\,d\mu$ and
$d_{j,q}=E(M_{j,q})-E(a_j)$. We have $z_{j,q}\uparrow g_j$ and
$0\le d_{j,q}\le R_j$. The exact identity
\[
 B-F_\mu(M_{j,q})=\varepsilon_j+z_{j,q}-d_{j,q}
\]
therefore implies
\[
 \varepsilon_j+g_j-R_j
 \le\sup_{q\ge j}[B-F_\mu(M_{j,q})]
 \le\varepsilon_j+g_j.
\]
This proves the first limit in \eqref{new:lostmax}. At fixed $j$,
$F_\mu(M_{j,q})\to L_j-\int S_j\,d\mu$; now let $j\to\infty$
to obtain the second. The equivalences follow, because each full
functional deficit is nonnegative. No positivity of $S_j$ was used.
\end{proof}

\section{Where an upper-relaxation excess must occur}
\begin{proposition}\label{new:spikes}
Let $R$ be the canonical ordinary USC relaxation of the maximizing
sequence, and let $x$ be full dimensional. If
$R(x)>\ell(x)+4\eta$ for some $\eta>0$, then for all sufficiently
large $m$ there is an ordinary neighborhood $W_m$ of $x$ such that,
for every $j\ge m$,
\[
 x\in\overline{\,O_m\cap W_m\cap
                  \{S_j>\ell(x)+4\eta\}\,}.
\]
The sets inside the closure are ordinary open and have capacity at
most $P_m$ on each fixed $\C_{C+k}$ once $m\ge k$.
This gives a precise necessary form of a remaining relaxation gap;
it proves neither that such a gap occurs nor that it cannot occur.
\end{proposition}
\begin{proof}
Since $x\notin N$, eventually $x\notin O_m$. Increase $m$ until
$|a_m(x)-\ell(x)|<\eta$ and $2^{1-m}<\eta$.
Continuity gives a neighborhood $W_m$ with
$a_m<\ell(x)+2\eta$. For $y\in W_m\setminus O_m$ and $n\ge m$,
the telescoping bound gives $a_n(y)<\ell(x)+3\eta$.
On the other hand the definition
\[
 R(x)=\inf_j\inf_{W\ni x}\sup_{y\in W,n\ge j}a_n(y)
\]
implies that every neighborhood of $x$ in $W_m$, and every $j\ge m$,
contains a $y$ with $a_n(y)>\ell(x)+4\eta$ for some $n\ge j$.
It must belong to $O_m$. This proves the closure assertion.
The threshold set of the LSC $S_j$ is open, and its intersection is
contained in $O_m$, so \eqref{new:capacity} proves the capacity bound.
Small capacity is not an ordinary neighborhood exclusion theorem.
\end{proof}

\section{A precise sufficient continuity bridge, with separate equation gate}
\begin{proposition}\label{new:bridge}
If the actual $S_j$ constructed above are continuous on all $X$, then
$S_j\in\C$, $E(S_j)=L_j$, and $S_j\downarrow u$ is an original
continuous decreasing-sequence presentation with finite lower energy.
Consequently $u\in\mathcal G$, $E(u)=e$, and the original canonical
generated MA measure is $\nu$. Its prescribed-measure objective is
$E(u)-\int u\,d\mu=B-J$.
The needed continuity is unproved. Even this sufficient analytic
representation input alone does not prove $\nu=\mu$.
\end{proposition}
\begin{proof}
At fixed $j$, the actual continuous $M_{j,q}$ increase to a continuous
$S_j$ on compact Hausdorff $X$. The increasing version of Dini's
argument is a finite-open-cover proof, and gives uniform convergence.
The entire minimal continuous CLD class is uniformly closed, so
$S_j\in\C$. CLD product continuity gives energy convergence and
$E(S_j)=L_j$. The sequence $S_j$ is decreasing, normalized and has
energies at least $-C$; thus its actual pointwise infimum belongs to
the original sequence-generated domain by its definition. Its energy
is $\lim L_j=e$. Passing the uniformly convergent bounded products
at fixed $j$ in \eqref{new:maximum} gives $I(a_j,S_j)\le R_j$.
The fixed-test Green estimate and density therefore imply
$\MA(S_j)\rightharpoonup\nu$ as $j\to\infty$. The original generated
operator is the canonical weak limit of a continuous decreasing
presentation, as reconstructed immediately below. Hence it is $\nu$.
The formula for its objective follows from Theorem \ref{new:Lone}.
No equation-identifying variation has been supplied.
\end{proof}

Here is a complete reconstruction of the particular generated-operator
fact used by this implication. For any continuous decreasing
presentation $c_n\downarrow v$ with finite lower energy, compact strict
domination compares a term of either presentation with a term of the
other plus an arbitrary positive constant. Monotonicity and
$E(c+\epsilon)=E(c)+2\epsilon$ prove presentation independence of
the limiting real energy. The same argument applies to finite convex
combinations. On their finite simplex the energies are polynomials of
degree at most three; convergence at a finite unisolvent set implies
coefficient convergence. Such a set exists because a polynomial
vanishing on an open simplex is identically zero.
For a fixed bounded $w\in\C$, set
$P_n(t)=E((1-t)c_n+tw)$ and $A_n(t)=E(tc_n)$.
Actual first variation gives
$\int w\MA(c_n)=P'_n(0)+A'_n(1)$.
Coefficient convergence makes this limit canonical. Density of
$\operatorname{span}\C$ and positive common mass two extend the limits
to every continuous test; compact-Hausdorff Riesz gives a canonical
actual positive Radon measure $\MA(v)$, exactly the weak limit of the
bounded analytic measures, of mass two. This construction applies to
an already actual decreasing presentation, and not to a merely Borel
or Banach representative.
For integrability, first pass bounded tangent inequalities with each
fixed continuous competitor: $D(\MA(v))\le E(v)-d$, where
$d=\lim\int c_n\MA(c_n)$ is the coefficient $A'_n(1)$.
After a common upper shift the bound is finite. Sequential monotone
convergence in the full dual inequality therefore makes $v$ integrable
against $\MA(v)$ and each other eligible measure.
For $q\ge n$, the ordered cocycle gives
$0\le\int(c_n-c_q)\MA(c_q)\le3(E(c_n)-E(c_q))$.
Pass first $q\to\infty$ against fixed continuous $c_n$, then
$n\to\infty$ using that integrability; this proves
$d=\int v\MA(v)$ and
$D(\MA(v))=E(v)-\int v\MA(v)$.
Every passage is sequential or against fixed continuous tests.
This reconstruction proves the operator fact used here, not the
historical comparison, full uniqueness, finite-extension, descent or
nonalgebraization claims, whose deeper audit remains unfinished.

The same bounded proofs independently give the canonical map
$\mathcal V(\mu)=\nu$ and its idempotence. The net maximizing for
$\mu$ also maximizes for $\nu$ by \eqref{dualcontinuous}; it already
has weak MA limit $\nu$, so maximizing-family independence gives
$\mathcal V(\nu)=\nu$. The exact one-sided cubic Taylor argument
above proves $D(\nu)\le D(\mu)$. An already generated measure
$\MA(v)$ is fixed: its decreasing presentation has limiting
objective $D(\MA(v))$ by the preceding paragraph, and its MA measures
converge to $\MA(v)$; normalization changes neither. Neither a fixed
point nor $J=0$ alone supplies the missing analytic or two-sided
variation inputs for the original equation.

\section{New scoped status}
The diagonal construction, quasi-uniform convergence, uniform integral
tail bounds, the exact $J$ defect and its finite-maximum objective
interpretation, and the relaxation-spike proposition are complete
deductions at the declared bounded analytic-foundation boundary.
The continuity bridge is an implication with an unproved antecedent.
No prescribed support has been reduced; $\mu$ is every positive
mass-two Radon measure with finite full bounded dual.
The Borel function $u$ is not asserted globally USC or CLD psh;
neither $S_j$ nor its ordinary USC regularization is inserted into an
unproved positivity class. The result does not replace analytic
weights by an abstract completion. We have not proved $J=0$, $K=0$,
$R=\ell$ on all full points, an adequate envelope theorem, or
$\MA(\text{actual analytic solution})=\mu$.
The genuinely nonalgebraizable and coherent-descent premise has not
been independently certified by these bounded deductions.
Original P16.7 and the complete deeper source audit remain open.
